\section{Method}
\label{sec:method}
\hans{Logic:
subsection 1
first introduce frame to frame ode flow
however, in forecasting, we always start from a history so in a ode flow we do not have any stochasticity to model uncertainty.
second introduce the vae with elbo
third, objective the crps loss also is a proper scoring rule.
subsection 2: training
subsection 3: inference
}

\subsection{Frame-to-frame ODE Flow}
\label{sec:frame-to-frame-method}
\hans{highlight learning on the physical-time flow instead of the}
From $N_{\mathrm{traj}}$ offline trajectories of lengths $T_n$, we form the
autonomous and controlled adjacent-frame datasets
\begin{align}
  \mathcal D_{\mathrm{ff}}
  &:={\left\{
      \bigl(x_k^{(n)},x_{k+1}^{(n)}\bigr)
      : n=1,\ldots,N_{\mathrm{traj}},\;
        k=1,\ldots,T_n-1
    \right\}},
  \label{eq:problem_uncontrolled_pairs}\\
  \mathcal D_{\mathrm{ff}}^{\mathrm{act}}
  &:={\left\{
      \bigl(x_k^{(n)},a_k^{(n)},x_{k+1}^{(n)}\bigr)
      : n=1,\ldots,N_{\mathrm{traj}},\;
        k=1,\ldots,T_n-1
    \right\}}.
  \label{eq:problem_controlled_pairs}
\end{align}
The trajectory identity and temporal order are retained when sampling local
transitions.  Following \citet{lim2026flow}, we construct the rollout from
these adjacent frames.
The discrete index $k$ denotes physical time, whereas
$\tau\in[0,1]$ parameterizes a single frame transition.  Each observed pair
defines the local interpolation
\begin{align}
  x_{k,\tau}
  &:=(1-\tau)x_k+\tau x_{k+1},
  \label{eq:method_frame_interpolation}\\
  u_{k,\tau}^{\mathrm{loc}}
  &:=\frac{\partial x_{k,\tau}}{\partial\tau}
   =x_{k+1}-x_k.
  \label{eq:method_frame_velocity}
\end{align}
This is a flow-matching path in the chosen state representation; it does not
assume that the unresolved physical motion between sampled frames is linear.

For consistency with Appendix~\ref{sec:clean_velocity_elbo}, consider a
training window $x_{a:T}$ with forecast origin $m$ and set $P:=T-m$.  Write
$k=m+r$ for $r=0,\ldots,P-1$, and define the global normalized coordinate
\begin{equation}
  s:=\frac{r+\tau}{P},
  \qquad
  s\in I_r:=\left[\frac rP,\frac{r+1}{P}\right).
  \label{eq:method_local_global_time}
\end{equation}
Then $x_s=x_{m+r,\tau}$, and the change of variables yields
\begin{equation}
  u_s:=\frac{dx_s}{ds}
  =P\frac{\partial x_{m+r,\tau}}{\partial\tau}
  =P(x_{m+r+1}-x_{m+r}).
  \label{eq:method_global_velocity}
\end{equation}
Equations~\ref{eq:method_local_global_time} and
\ref{eq:method_global_velocity} are exactly the construction used in
Eqs.~\ref{eq:clean_interpolation}--\ref{eq:clean_velocity}.  Thus $s$ is the
normalized trajectory coordinate and is never used as a system-state symbol.

Let $h_k:=H(x_{a:k})$ be a deterministic representation of the available
$x$-history.  It is model context, not another physical state or a Bayesian
belief.  Conditional on the trajectory-level generative variable $z$
introduced below, the autonomous velocity defines the frame update
\begin{equation}
  \frac{d\widehat x_s}{ds}
  =v_\theta(\widehat x_s,s,h_m,z),
  \qquad
  s\in I_r,
  \qquad
  \left.\widehat x_s\right|_{s=r/P}=\widehat x_{m+r},
  \label{eq:method_frame_flow_ode}
\end{equation}
and emits
\begin{equation}
  \widehat x_{m+r+1}
  :=\left.\widehat x_s\right|_{s=(r+1)/P}
  =\widehat x_{m+r}
   +\int_{r/P}^{(r+1)/P}
      v_\theta(\widehat x_s,s,h_m,z)\,ds.
  \label{eq:method_frame_flow_map}
\end{equation}
Following the action-conditioned format of \citet{maes2026leworldmodel}, the
controlled model supplies the aligned action explicitly to the transition.
Accordingly, its right-hand side is
$v_\theta(\widehat x_s,s,h_m,a_{m+r},z)$; the action is not
absorbed into the history encoder $H$.
The observed context $h_m=H(x_{a:m})$ is fixed throughout the forecast, in
agreement with Appendix~\ref{sec:clean_velocity_elbo}; generated frames enter
the ODE through $\widehat x_s$ rather than through a refreshed context
encoding.  Repeating Eq.~\ref{eq:method_frame_flow_map} for
$r=0,\ldots,P-1$ produces one trajectory sample.

\subsection{Variational ODE Flow}
\hans{While learning from a frame to frame flow matching allow to }
Let $d=\dim(u_s)$ and $p(z)=\mathcal N(0,I_{d_z})$.  The training posterior is
\begin{equation}
  q_\phi(z\mid u_s,x_s,s,h)
  =\mathcal N\!\left(
    z;
    \mu_\phi(u_s,x_s,s,h),
    \operatorname{diag}\!\bigl(\sigma_\phi^2(u_s,x_s,s,h)\bigr)
  \right).
  \label{eq:method-variational-posterior}
\end{equation}

The nonlinear decoder parameterizes the mean of the fixed-variance conditional
velocity likelihood
\begin{equation}
  p_\theta(u_s\mid x_s,s,h,z)
  =\mathcal N\!\left(
    u_s;
    v_\theta(x_s,s,h,z),
    \sigma_u^2 I_d
  \right),
  \label{eq:method-velocity-likelihood}
\end{equation}
where $\sigma_u^2>0$ is fixed.  Thus, conditional on $z$, $v_\theta$
deterministically predicts the likelihood mean.  Marginalizing the random
latent gives a Gaussian location mixture, which can be non-Gaussian and
multimodal.
For controlled trajectories, the aligned $a_{m+r}$ is an additional,
observed conditioning variable in both networks:
\begin{equation}
  q_\phi(z\mid u_s,x_s,s,h,a_{m+r}),
  \qquad
  p_\theta(u_s\mid x_s,s,h,a_{m+r},z),
  \qquad s\in I_r.
  \label{eq:method-controlled-conditionals}
\end{equation}
All autonomous equations are recovered by omitting this action argument; the
ELBO proof applies conditionally on the supplied action.


\paragraph{Generalized Energy Loss}
\hans{Motivation: }
\begin{equation}
  \widehat{\mathcal L}_{\beta,M}^{q}
  :=\frac1M\sum_{i=1}^M
       \left|\widehat U_s^{(i)}-u_s\right|^\beta
    -\frac{\lambda}{M^2}\sum_{i,j=1}^M
       \left|\widehat U_s^{(i)}-\widehat U_s^{(j)}\right|^\beta.
  \label{eq:generalized_crps_training_loss}
\end{equation}


Algorithms~\ref{alg:vsf-training} and~\ref{alg:vsf-inference} summarize
training and a $P$-step parallel forecast, respectively.  Here $P$ is the
number of future sampling steps, $P=T-m$,
and ensemble members are indexed by $j$ because $m$ denotes the forecast
origin.  The inference panel shows the knot-aligned Euler specialization that
corresponds directly to the proof; additional numerical substeps may be used
without inserting them into the observation history.

% Lightweight ruled pseudocode, kept local to this section so that no
% additional algorithm package is required.
\newcounter{vsfalgorithm}
\newcommand{\vsfalgorithmheader}[2]{%
  \refstepcounter{vsfalgorithm}\label{#1}%
  \par\noindent\textbf{Algorithm \thevsfalgorithm: #2}\par
}
\newenvironment{vsfalgorithmsteps}{%
  \begin{list}{}{
    \setlength{\leftmargin}{1.0em}%
    \setlength{\rightmargin}{0pt}%
    \setlength{\labelwidth}{0pt}%
    \setlength{\labelsep}{0pt}%
    \setlength{\itemsep}{0.15em}%
    \setlength{\parsep}{0pt}%
    \setlength{\topsep}{0.15em}%
    \setlength{\partopsep}{0pt}%
  }%
}{\end{list}}

\begin{figure}[t]
  \centering
  \begin{minipage}[t]{0.485\linewidth}
    \vspace{0pt}
    \footnotesize
    \hrule height 0.8pt
    \vspace{0.25em}
    \vsfalgorithmheader{alg:vsf-training}{Training}
    \vspace{0.15em}\hrule height 0.4pt\vspace{0.3em}
    \noindent\textbf{Input:} Windows $x_{a:T}\sim p_{\mathrm{data}}$ (paired
    with $a_{a:T-1}$ for controlled data), prior $p(z)$, and
    $\beta_{\mathrm{KL}}$.

    \noindent\textbf{Output:} Trained $q_\phi$ and $v_\theta$.

    \noindent\textbf{while} not converged \textbf{do}
    \begin{vsfalgorithmsteps}
      \item[] Sample $x_{a:T}$ and $s\sim\mathrm{Unif}[0,1]$; set
      $h\leftarrow H(x_{a:m})$.

      \item[] Evaluate $(x_s,u_s)$ on the clean trajectory, where
      $u_s=\mathrm{d}x_s/\mathrm{d}s$.

      \item[] Set the conditioning tuple
      \[
        c_s\leftarrow
        \begin{cases}
          h, & \text{autonomous},\\
          (h,a_{m+r}), & \text{controlled, with }s\in I_r.
        \end{cases}
      \]

      \item[] $(\mu_\phi,\sigma_\phi)\leftarrow
      \operatorname{Enc}_\phi(u_s,x_s,s,c_s)$ and set
      $q_{\phi,s}:=q_\phi(\,\cdot\mid u_s,x_s,s,c_s)$.

      \item[] Sample $\eta\sim\mathcal N(0,I_{d_z})$ and set
      $z\leftarrow\mu_\phi+\sigma_\phi\odot\eta$.

      \item[] Predict
      $\widehat u_s\leftarrow v_\theta(x_s,s,c_s,z)$.

      \item[] Compute
      \[
        \mathcal L\leftarrow
        \frac{1}{d}\|\widehat u_s-u_s\|_2^2
        +\beta_{\mathrm{KL}}\KL(q_{\phi,s}\|p).
      \]

      \item[] Update $(\theta,\phi)$ using $\nabla_{\theta,\phi}\mathcal L$.
    \end{vsfalgorithmsteps}
    \noindent\textbf{end while}
    \vspace{0.3em}\hrule height 0.8pt
  \end{minipage}
  \hfill
  \begin{minipage}[t]{0.485\linewidth}
    \vspace{0pt}
    \footnotesize
    \hrule height 0.8pt
    \vspace{0.25em}
    \vsfalgorithmheader{alg:vsf-inference}{Inference}
    \vspace{0.15em}\hrule height 0.4pt\vspace{0.3em}
    \noindent\textbf{Input:} Trained $v_\theta$, observed context
    $x_{a:m}$, prior $p(z)$, horizon $P=T-m$, ensemble size $M$, and supplied
    $a_{m:T-1}$ for controlled inference.

    \noindent\textbf{Output:} Forecast ensemble
    $\{\widehat x_{m+1:T}^{(j)}\}_{j=1}^{M}$.

    \begin{vsfalgorithmsteps}
      \item[] Set $h\leftarrow H(x_{a:m})$.

      \item[] For $j=1,\ldots,M$, sample $z^{(j)}\sim p(z)$ once and set
      \[
        \widehat x_m^{(j)}\leftarrow x_m.
      \]

      \item[] \textbf{for} $k=0,\ldots,P-1$ \textbf{do}
      \begin{vsfalgorithmsteps}
        \item[] Set $s_k\leftarrow k/P$ and $\Delta s\leftarrow1/P$.

        \item[] In parallel for $j=1,\ldots,M$, evaluate
        \[
          \widehat u_{s_k}^{(j)}\leftarrow
          v_\theta\!\left(
            \widehat x_{m+k}^{(j)},s_k,
            h,z^{(j)}
          \right).
        \]

        \item[] For controlled inference, instead set
        $\widehat u_{s_k}^{(j)}\leftarrow
         v_\theta(\widehat x_{m+k}^{(j)},s_k,h,a_{m+k},z^{(j)})$.

        \item[] Advance and emit
        \[
          \widehat x_{m+k+1}^{(j)}\leftarrow
          \widehat x_{m+k}^{(j)}
          +\Delta s\,\widehat u_{s_k}^{(j)}.
        \]

      \end{vsfalgorithmsteps}
      \item[] \textbf{end for}
    \end{vsfalgorithmsteps}
    \noindent Return the $M$ trajectories; keep each $z^{(j)}$ fixed over
    its complete trajectory.
    \vspace{0.3em}\hrule height 0.8pt
  \end{minipage}
\end{figure}

\hans{add crps loss}


\
% The fixed context h_m matches Appendix~\ref{sec:clean_velocity_elbo}.  For an
% open-ended stream, advance the forecast origin and recompute h_m from the newly
% observed window; a new latent is drawn when a new ensemble forecast is
% initialized.
